% questions.tex —— 1991 年·全国卷·理科 · 题目内容片段（无答案）
% 由 tools/generate_projects.py 依来源总账自 1991/1991_数学真题_全国理_试题.md 生成，勿手改（改源 MD 后重跑）


\gkpart{试题}

\begin{gktip}
\textbf{来源说明}：固定来源为 \texttt{普通高考/1991/1991全国理.pdf}，通过题面元数据中的来源编号关联来源总账。原件为一页 A3 横向排印稿，印刷页码 39，含第 1—26 题；原件未附答案、解析或分值。以下题面按原页转录，答案和解析另行标为非官方推导。

\textbf{图形说明}：原页未印题图；第 23 题按文字条件作答，不从示意图推断未给出的度量条件。
\end{gktip}

\begin{question}
\gkti{1}

已知 $\sin\alpha=\dfrac45$，并且 $\alpha$ 是第二象限的角，那么 $\tan\alpha$ 的值等于\mbox{（　　）}

（A）$-\dfrac43$　　（B）$-\dfrac34$　　（C）$\dfrac34$　　（D）$\dfrac43$

\end{question}
\begin{question}
\gkti{2}

焦点在 $(-1,0)$，顶点在 $(1,0)$ 的抛物线方程是\mbox{（　　）}

（A）$y^2=8(x+1)$　　（B）$y^2=-8(x+1)$

（C）$y^2=8(x-1)$　　（D）$y^2=-8(x-1)$

\end{question}
\begin{question}
\gkti{3}

函数 $y=\cos^4x-\sin^4x$ 的最小正周期是\mbox{（　　）}

（A）$\dfrac\pi2$　　（B）$\pi$　　（C）$2\pi$　　（D）$4\pi$

\end{question}
\begin{question}
\gkti{4}

如果把两条异面直线看成“一对”，那么六棱锥的棱所在的 12 条直线中，异面直线共有\mbox{（　　）}

（A）12 对　　（B）24 对　　（C）36 对　　（D）48 对

\end{question}
\begin{question}
\gkti{5}

函数 $y=\sin\left(2x+\dfrac{5\pi}{2}\right)$ 的图象的一条对称轴的方程是\mbox{（　　）}

（A）$x=-\dfrac\pi2$　　（B）$x=-\dfrac\pi4$　　（C）$x=\dfrac\pi8$　　（D）$x=\dfrac{5\pi}4$

\end{question}
\begin{question}
\gkti{6}

如果三棱锥 $S-ABC$ 的底面是不等边三角形，侧面与底面所成的二面角都相等，且顶点 $S$ 在底面的射影 $O$ 在 $\triangle ABC$ 内，那么 $O$ 是 $\triangle ABC$ 的\mbox{（　　）}

（A）垂心　　（B）重心　　（C）外心　　（D）内心

\end{question}
\begin{question}
\gkti{7}

已知 $\{a_n\}$ 是等比数列，且 $a_n>0$，$a_2a_4+2a_3a_5+a_4a_6=25$，那么 $a_3+a_5$ 的值等于\mbox{（　　）}

（A）5　　（B）10　　（C）15　　（D）20

\end{question}
\begin{question}
\gkti{8}

如果圆锥曲线的极坐标方程为 $\rho=\dfrac{16}{5-3\cos\theta}$，那么它的焦点的极坐标为\mbox{（　　）}

（A）$(0,0)$、$(6,\pi)$　　（B）$(-3,0)$、$(3,0)$

（C）$(0,0)$、$(3,0)$　　（D）$(0,0)$、$(6,0)$

\end{question}
\begin{question}
\gkti{9}

从 4 台甲型和 5 台乙型电视机中任意取出 3 台，其中至少要有甲型与乙型电视机各 1 台，则不同的取法共有\mbox{（　　）}

（A）140 种　　（B）84 种　　（C）70 种　　（D）35 种

\end{question}
\begin{question}
\gkti{10}

如果 $AC<0$ 且 $BC<0$，那么直线 $Ax+By+C=0$ 不通过\mbox{（　　）}

（A）第一象限　　（B）第二象限　　（C）第三象限　　（D）第四象限

\end{question}
\begin{question}
\gkti{11}

设甲、乙、丙是三个命题。如果甲是乙的必要条件；丙是乙的充分条件但不是乙的必要条件，那么\mbox{（　　）}

（A）丙是甲的充分条件，但不是甲的必要条件

（B）丙是甲的必要条件，但不是甲的充分条件

（C）丙是甲的充要条件

（D）丙不是甲的充分条件，也不是甲的必要条件

\end{question}
\begin{question}
\gkti{12}

\nopagebreak[4]
\[
\lim_{n\to\infty}\left[n\left(1-\frac13\right)\left(1-\frac14\right)\left(1-\frac15\right)\cdots\left(1-\frac1{n+2}\right)\right]
\]

的值等于\mbox{（　　）}

（A）0　　（B）1　　（C）2　　（D）3

\end{question}
\begin{question}
\gkti{13}

如果奇函数 $f(x)$ 在区间 $[3,7]$ 上是增函数且最小值为 5，那么 $f(x)$ 在区间 $[-7,-3]$ 上是\mbox{（　　）}

（A）增函数且最小值为 $-5$　　（B）增函数且最大值为 $-5$

（C）减函数且最小值为 $-5$　　（D）减函数且最大值为 $-5$

\end{question}
\begin{question}
\gkti{14}

圆 $x^2+2x+y^2+4y-3=0$ 上到直线 $x+y+1=0$ 的距离为 $\sqrt2$ 的点共有\mbox{（　　）}

（A）1 个　　（B）2 个　　（C）3 个　　（D）4 个

\end{question}
\begin{question}
\gkti{15}

设全集为 $\mathbb R$，$f(x)=\sin x$，$g(x)=\cos x$，$M=\{x\mid f(x)\ne0\}$，$N=\{x\mid g(x)\ne0\}$，那么集合 $\{x\mid f(x)g(x)=0\}$ 等于\mbox{（　　）}

（A）$\overline M\cap\overline N$　　（B）$\overline M\cup N$　　（C）$M\cup\overline N$　　（D）$\overline M\cup\overline N$

\end{question}
\begin{question}
\gkti{16}

$\arctan\dfrac13+\arctan\dfrac12$ 的值是\underline{\hspace{6em}}。

\end{question}
\begin{question}
\gkti{17}

不等式 $6^{x^2+x-2}<1$ 的解集是\underline{\hspace{8em}}。

\end{question}
\begin{question}
\gkti{18}

已知正三棱台上底面边长为 2，下底面边长为 4，且侧棱与底面所成的角是 $45^\circ$，那么这个正三棱台的体积等于\underline{\hspace{6em}}。

\end{question}
\begin{question}
\gkti{19}

$(ax+1)^7$ 的展开式中，$x^3$ 的系数是 $x^2$ 的系数与 $x^4$ 的系数的等差中项。若实数 $a>1$，那么 $a=$\underline{\hspace{6em}}。

\end{question}
\begin{question}
\gkti{20}

在球面上有四个点 $P$、$A$、$B$、$C$，如果 $PA$、$PB$、$PC$ 两两互相垂直，且 $PA=PB=PC=a$，那么这个球面的面积是\underline{\hspace{6em}}。

\end{question}
\begin{question}
\gkti{21}

求函数 $y=\sin^2x+2\sin x\cos x+3\cos^2x$ 的最小值，并写出使函数 $y$ 取最小值的 $x$ 的集合。

\end{question}
\begin{question}
\gkti{22}

已知复数 $z=1+i$，求复数 $\dfrac{z^2-3z+6}{z+1}$ 的模和辐角的主值。

\end{question}
\begin{question}
\gkti{23}

已知 $ABCD$ 是边长为 4 的正方形，$E$、$F$ 分别是 $AB$、$AD$ 的中点，$GC$ 垂直于 $ABCD$ 所在的平面，且 $GC=2$。求点 $B$ 到平面 $EFG$ 的距离。

\end{question}
\begin{question}
\gkti{24}

根据函数单调性的定义，证明函数 $f(x)=-x^3+1$ 在 $(-\infty,+\infty)$ 上是减函数。

\end{question}
\begin{question}
\gkti{25}

\Needspace{5\baselineskip}
已知 $n$ 为自然数，实数 $a>1$，解关于 $x$ 的不等式
\nopagebreak[4]
\[
\log_a x-4\log_{a^2}x+12\log_{a^3}x+\cdots+n(-2)^{n-1}\log_{a^n}x
>
\frac{1-(-2)^n}{3}\log_a(x^2-a)
\]

\end{question}
\begin{question}
\gkti{26}

双曲线的中心在坐标原点 $O$，焦点在 $x$ 轴上，过双曲线右焦点且斜率为 $\sqrt{\dfrac35}$ 的直线交双曲线于 $P$、$Q$ 两点。若 $OP\perp OQ$，$|PQ|=4$，求双曲线的方程。

\end{question}